\chapter{Full-Length University Mock Examinations}
\label{chap:mock-examinations}

% ==============================================================================
% SECTION 9.1: EXAMINATION STRUCTURE
% ==============================================================================
\section{University Examination Structure \& Instructions}

This chapter provides six complete, realistic 100-mark mock examinations modeled after the Lovely Professional University (LPU) Semester 1 final examination paper pattern. Each examination is followed by complete, step-by-step worked scoring solutions.

\begin{conceptbox}[Standard Examination Instructions]
\begin{itemize}[leftmargin=1.5em]
    \item \textbf{Duration}: 3 Hours $\cdot$ \textbf{Maximum Marks}: 100 per paper.
    \item \textbf{Section A}: 10 Short Compulsory Questions ($10 \times 2 = 20$ Marks).
    \item \textbf{Section B}: 6 Analytical Medium Problems ($6 \times 5 = 30$ Marks).
    \item \textbf{Section C}: 5 Long Comprehensive Programming Problems ($5 \times 10 = 50$ Marks).
\end{itemize}
\end{conceptbox}

% ==============================================================================
% MOCK EXAMINATION 1
% ==============================================================================
\section{University Mock Examination 1}

\begin{mockexam}[INT108 Mock Examination 1 (100 Marks --- 3 Hours)]
\noindent\textbf{Section A: Short Compulsory Questions ($10 \times 2 = 20\text{ Marks}$)}
\begin{enumerate}[leftmargin=1.5em]
    \item State the difference between integer division in Python 2 and Python 3.
    \item What is the purpose of adding Python to the system PATH environment variable?
    \item Differentiate between the \texttt{range(stop)} and \texttt{random.randint(a, b)} upper boundary rules.
    \item Explain the output of \texttt{bool("")} and \texttt{bool("0")}.
    \item What is a void function, and what value does it implicitly return in Python?
    \item Explain the LEGB scope resolution rule.
    \item Why cannot a list be used as a dictionary key?
    \item Differentiate between \texttt{pop()} and \texttt{remove()} methods of a list.
    \item What is Name Mangling in Python classes, and why is it used?
    \item Differentiate between \texttt{re.match()} and \texttt{re.search()}.
\end{enumerate}

\medskip
\noindent\textbf{Section B: Analytical Medium Problems ($6 \times 5 = 30\text{ Marks}$)}
\begin{enumerate}[leftmargin=1.5em, start=11]
    \item Predict the exact output of:
    \begin{lstlisting}[language=Python3]
a = [1, 2, 3]
b = a
c = a[:]
b[0] = 99
c[1] = 88
print(a, b, c)
    \end{lstlisting}
    \item Write a function \texttt{is\_prime(n)} that returns \texttt{True} if $n$ is prime, using a \texttt{for...else} loop construct.
    \item Explain why Python does not support traditional method overloading and show how to simulate it using default arguments.
    \item Write a Python function using \texttt{try-except-else-finally} to safely divide two input strings as floating-point numbers.
    \item Explain the output of: \texttt{print(2 ** 3 ** 2, 16 // 5, -16 // 5)}.
    \item Write a regular expression to extract all URLs starting with \texttt{http://} or \texttt{https://} from a text string.
\end{enumerate}

\medskip
\noindent\textbf{Section C: Comprehensive Programming Problems ($5 \times 10 = 50\text{ Marks}$)}
\begin{enumerate}[leftmargin=1.5em, start=17]
    \item Write a recursive function \texttt{fibonacci(n)} to return the $n$-th Fibonacci number, and trace the complete call stack for $n=4$. (10 Marks)
    \item Design a complete \texttt{Matrix} class with methods to initialize an $N \times M$ matrix, add two matrices, and compute the transpose. (10 Marks)
    \item Create an inheritance hierarchy: Base class \texttt{Account} with private \texttt{\_\_balance}, and derived class \texttt{SavingsAccount} with an added \texttt{interest\_rate} and \texttt{add\_interest()} method. (10 Marks)
    \item Write a program that reads a text file named \texttt{source.txt}, counts the occurrences of every unique word using a dictionary, and writes the sorted word counts to \texttt{word\_summary.txt}. (10 Marks)
    \item Write a Python script to validate user registration data using regular expressions:
    \begin{enumerate}[label=(\alph*)]
        \item Email address validation (e.g., \texttt{user@domain.com}). (5 Marks)
        \item 10-digit mobile number starting with 6, 7, 8, or 9. (5 Marks)
    \end{enumerate}
\end{enumerate}
\end{mockexam}

\subsection{Complete Worked Solutions for Mock Examination 1}

\begin{solutionbox}[Solutions to Mock Exam 1 (100 Marks)]
\noindent\textbf{Section A Solutions ($10 \times 2 = 20\text{ Marks}$)}:
\begin{enumerate}
    \item In Python 2, \texttt{5 / 2} truncated to integer \texttt{2}. In Python 3, \texttt{5 / 2} returns float \texttt{2.5}. [2M]
    \item Adding to PATH allows the CLI to recognize \texttt{python} from any directory. [2M]
    \item \texttt{range(1, 5)} excludes 5; \texttt{randint(1, 5)} includes both 1 and 5. [2M]
    \item \texttt{bool("")} is \texttt{False} (empty string); \texttt{bool("0")} is \texttt{True} (non-empty string). [2M]
    \item A void function performs tasks without returning a value; it implicitly returns \texttt{None}. [2M]
    \item LEGB: Local $\to$ Enclosing $\to$ Global $\to$ Built-in. [2M]
    \item Lists are mutable and unhashable; dictionary keys must be hashable. [2M]
    \item \texttt{pop(i)} removes by index and returns it; \texttt{remove(val)} removes first occurrence of value. [2M]
    \item Double underscore attributes (\texttt{\_\_x}) are mangled to \texttt{\_Class\_\_x} to prevent accidental override. [2M]
    \item \texttt{match()} checks strictly at start (index 0); \texttt{search()} scans entire string. [2M]
\end{enumerate}

\medskip
\noindent\textbf{Section B Solutions ($6 \times 5 = 30\text{ Marks}$)}:
\begin{enumerate}[start=11]
    \item \texttt{b=a} is an alias; \texttt{c=a[:]} is a shallow copy. Output: \texttt{[99, 2, 3] [99, 2, 3] [1, 88, 3]}. [5M]
    \item Prime function using \texttt{for...else}:
    \begin{lstlisting}[language=Python3]
def is_prime(n):
    if n <= 1: return False
    for i in range(2, int(n**0.5) + 1):
        if n % i == 0:
            return False
    else:
        return True
    \end{lstlisting} [5M]
    \item Python is dynamically typed; subsequent definitions overwrite earlier ones. Simulated via default parameters (\texttt{def add(a, b=0): return a + b}). [5M]
    \item Safe division:
    \begin{lstlisting}[language=Python3]
def safe_div(s1, s2):
    try:
        a, b = float(s1), float(s2)
        res = a / b
    except ZeroDivisionError: print("Division by Zero")
    except ValueError: print("Invalid Number")
    else: print(f"Result: {res}")
    finally: print("Done")
    \end{lstlisting} [5M]
    \item $2^{(3^2)} = 2^9 = \mathbf{512}$, $16 // 5 = \mathbf{3}$, $-16 // 5 = \mathbf{-4}$. [5M]
    \item URL regex: \texttt{r"https?://[\textbackslash w\textbackslash.-]+(?:/[\textbackslash w\textbackslash.-]*)*"}. [5M]
\end{enumerate}

\medskip
\noindent\textbf{Section C Solutions ($5 \times 10 = 50\text{ Marks}$)}:
\begin{enumerate}[start=17]
    \item \textbf{Fibonacci \& Stack}: $F(4) = F(3) + F(2) = (F(2)+F(1)) + (F(1)+F(0)) = (1+1) + (1+0) = \mathbf{3}$. Stack pushes $F(4) \to F(3) \to F(2) \to F(1) \to F(0)$ and unwinds sums. [10M]
    \item \textbf{Matrix Class}:
    \begin{lstlisting}[language=Python3]
class Matrix:
    def __init__(self, data):
        self.data = data
        self.rows = len(data)
        self.cols = len(data[0])
    def add(self, other):
        return Matrix([[self.data[r][c] + other.data[r][c] 
                        for c in range(self.cols)] for r in range(self.rows)])
    def transpose(self):
        return Matrix([[self.data[r][c] for r in range(self.rows)] 
                        for c in range(self.cols)])
    \end{lstlisting} [10M]
    \item \textbf{Account Inheritance}:
    \begin{lstlisting}[language=Python3]
class Account:
    def __init__(self, acc_id, balance):
        self.acc_id = acc_id
        self.__balance = float(balance)
    def get_balance(self): return self.__balance
    def deposit(self, amt): self.__balance += amt

class SavingsAccount(Account):
    def __init__(self, acc_id, balance, rate=0.04):
        super().__init__(acc_id, balance)
        self.rate = rate
    def add_interest(self):
        interest = self.get_balance() * self.rate
        self.deposit(interest)
    \end{lstlisting} [10M]
    \item \textbf{Word Count File Processing}:
    \begin{lstlisting}[language=Python3]
freq = {}
with open("source.txt", "r") as f:
    for line in f:
        for word in line.lower().split():
            clean = word.strip(".,!?:;\"'")
            freq[clean] = freq.get(clean, 0) + 1
with open("word_summary.txt", "w") as out:
    for w, c in sorted(freq.items()):
        out.write(f"{w}: {c}\n")
    \end{lstlisting} [10M]
    \item \textbf{Regex Validator}:
    \begin{lstlisting}[language=Python3]
import re
def validate_email(email):
    return bool(re.match(r"^[\w\.-]+@[\w\.-]+\.[a-zA-Z]{2,}$", email))

def validate_phone(phone):
    return bool(re.match(r"^[6-9]\d{9}$", phone))
    \end{lstlisting} [10M]
\end{enumerate}
\end{solutionbox}

% ==============================================================================
% MOCK EXAMINATION 2
% ==============================================================================
\section{University Mock Examination 2}

\begin{mockexam}[INT108 Mock Examination 2 (100 Marks --- 3 Hours)]
\noindent\textbf{Section A: Short Compulsory Questions ($10 \times 2 = 20\text{ Marks}$)}
\begin{enumerate}[leftmargin=1.5em]
    \item Explain the term Dynamic Typing in Python.
    \item Write the output of \texttt{print("Spam" * 3 + "Eggs")}.
    \item State the difference between \texttt{break} and \texttt{continue}.
    \item What is the difference between shallow copy and deep copy?
    \item State the purpose of the \texttt{math.ceil()} function.
    \item Write the syntax for creating a single-element tuple containing the number 99.
    \item What is the output of \texttt{dict.get("key", "default")} if \texttt{"key"} is present?
    \item Differentiate between \texttt{class} attributes and \texttt{instance} attributes.
    \item What does \texttt{pickle.dump()} do?
    \item Explain the meaning of the regex metacharacter \texttt{+} vs \texttt{*}.
\end{enumerate}

\medskip
\noindent\textbf{Section B: Analytical Medium Problems ($6 \times 5 = 30\text{ Marks}$)}
\begin{enumerate}[leftmargin=1.5em, start=11]
    \item Predict the output:
    \begin{lstlisting}[language=Python3]
def func(x, l=[]):
    l.append(x)
    return l
print(func(1))
print(func(2))
    \end{lstlisting}
    \item Write a Python program to extract all prime numbers in a list using a list comprehension.
    \item Implement a function \texttt{count\_vowels\_consonants(text)} that returns a tuple of counts.
    \item Write a function to check if a matrix is symmetric ($A = A^T$).
    \item Explain how the \texttt{with} context manager manages file closing under exceptions.
    \item Write a regex expression to extract all hashtags (\texttt{\#topic}) from a social media post.
\end{enumerate}

\medskip
\noindent\textbf{Section C: Comprehensive Programming Problems ($5 \times 10 = 50\text{ Marks}$)}
\begin{enumerate}[leftmargin=1.5em, start=17]
    \item Write a recursive function to solve the Greatest Common Divisor ($\gcd$) of two numbers, tracing $\gcd(48, 18)$. (10 Marks)
    \item Design a complete \texttt{ComplexNumber} class with \texttt{\_\_init\_\_}, \texttt{add()}, \texttt{multiply()}, and \texttt{\_\_str\_\_} methods. (10 Marks)
    \item Implement a complete \texttt{SparseMatrix} representation using coordinate dictionaries with an \texttt{add()} method. (10 Marks)
    \item Write a program that reads a CSV file containing sales data (\texttt{Item,Quantity,Price}), computes total revenue per item, and writes a summary report. (10 Marks)
    \item Build a complete automated quiz engine: store questions, options, and correct answers in a list of dictionaries, shuffle using \texttt{random}, administer quiz, and report final percentage score. (10 Marks)
\end{enumerate}
\end{mockexam}

\subsection{Complete Worked Solutions for Mock Examination 2}

\begin{solutionbox}[Solutions to Mock Exam 2 (100 Marks)]
\noindent\textbf{Section A Solutions ($10 \times 2 = 20\text{ Marks}$)}:
\begin{enumerate}
    \item Dynamic typing means variable types are checked and bound at runtime without explicit type declarations. [2M]
    \item \texttt{"SpamSpamSpamEggs"}. [2M]
    \item \texttt{break} exits the loop immediately; \texttt{continue} skips to the next iteration. [2M]
    \item Shallow copy clones outer container but references inner objects; deep copy clones all nested objects recursively. [2M]
    \item \texttt{ceil(x)} returns the smallest integer $\ge x$. [2M]
    \item \texttt{t = (99,)}. [2M]
    \item Returns the stored value associated with \texttt{"key"}. [2M]
    \item Class attributes are shared across all instances; instance attributes are unique to each object. [2M]
    \item Serializes a Python object into a binary byte stream and writes to disk. [2M]
    \item \texttt{+} matches 1 or more repetitions; \texttt{*} matches 0 or more repetitions. [2M]
\end{enumerate}

\medskip
\noindent\textbf{Section B Solutions ($6 \times 5 = 30\text{ Marks}$)}:
\begin{enumerate}[start=11]
    \item Mutable default argument trap: list \texttt{l} persists across calls. Output: \texttt{[1]}, then \texttt{[1, 2]}. [5M]
    \item Prime list: \texttt{[x for x in nums if x > 1 and all(x \% i != 0 for i in range(2, int(x**0.5)+1))]}. [5M]
    \item Vowel counter:
    \begin{lstlisting}[language=Python3]
def count_vowels_consonants(text):
    v = sum(1 for c in text.lower() if c in "aeiou")
    c = sum(1 for c in text.lower() if c.isalpha() and c not in "aeiou")
    return v, c
    \end{lstlisting} [5M]
    \item Symmetric matrix check: \texttt{all(matrix[r][c] == matrix[c][r] for r in range(n) for c in range(n))}. [5M]
    \item \texttt{with} invokes \texttt{\_\_enter\_\_} on open and guarantees \texttt{\_\_exit\_\_} runs upon exiting the block, automatically calling \texttt{close()}. [5M]
    \item Hashtag regex: \texttt{r"\#\textbackslash w+"}. [5M]
\end{enumerate}

\medskip
\noindent\textbf{Section C Solutions ($5 \times 10 = 50\text{ Marks}$)}:
\begin{enumerate}[start=17]
    \item \textbf{Recursive GCD}:
    \begin{lstlisting}[language=Python3]
def gcd(a, b):
    if b == 0:
        return a
    return gcd(b, a % b)
    \end{lstlisting}
    Trace: $\gcd(48, 18) \to \gcd(18, 12) \to \gcd(12, 6) \to \gcd(6, 0) = \mathbf{6}$. [10M]
    \item \textbf{ComplexNumber Class}:
    \begin{lstlisting}[language=Python3]
class ComplexNumber:
    def __init__(self, real, imag):
        self.real = real
        self.imag = imag
    def add(self, o):
        return ComplexNumber(self.real + o.real, self.imag + o.imag)
    def multiply(self, o):
        return ComplexNumber(self.real*o.real - self.imag*o.imag,
                             self.real*o.imag + self.imag*o.real)
    def __str__(self):
        return f"{self.real} + {self.imag}j"
    \end{lstlisting} [10M]
    \item \textbf{Sparse Matrix Dictionary Class}:
    \begin{lstlisting}[language=Python3]
class SparseMatrix:
    def __init__(self, rows, cols):
        self.rows, self.cols = rows, cols
        self.data = {}
    def set(self, r, c, val):
        if val != 0: self.data[(r, c)] = val
    def get(self, r, c): return self.data.get((r, c), 0)
    def add(self, o):
        res = SparseMatrix(self.rows, self.cols)
        all_keys = set(self.data.keys()).union(o.data.keys())
        for k in all_keys:
            res.set(k[0], k[1], self.get(k[0], k[1]) + o.get(k[0], k[1]))
        return res
    \end{lstlisting} [10M]
    \item \textbf{Sales CSV Processor}:
    \begin{lstlisting}[language=Python3]
totals = {}
with open("sales.csv", "r") as f:
    for line in f:
        parts = line.strip().split(",")
        if len(parts) == 3 and parts[1].isdigit():
            item, qty, price = parts[0], int(parts[1]), float(parts[2])
            totals[item] = totals.get(item, 0.0) + (qty * price)
with open("sales_summary.txt", "w") as out:
    for item, rev in totals.items():
        out.write(f"{item}: Rs. {rev:.2f}\n")
    \end{lstlisting} [10M]
    \item \textbf{Automated Quiz Engine}:
    \begin{lstlisting}[language=Python3]
import random
quiz = [
    {"q": "Keyword for function?", "opts": ["func", "def", "fn"], "ans": "def"},
    {"q": "Type of 10/2 in Py3?", "opts": ["int", "float", "num"], "ans": "float"}
]
random.shuffle(quiz)
score = 0
for item in quiz:
    print("\n" + item["q"])
    for i, opt in enumerate(item["opts"], 1): print(f"{i}. {opt}")
    choice = input("Your answer: ").strip()
    if choice.lower() == item["ans"].lower():
        score += 1
print(f"\nFinal Score: {score}/{len(quiz)} ({score/len(quiz)*100:.1f}%)")
    \end{lstlisting} [10M]
\end{enumerate}
\end{solutionbox}

% ==============================================================================
% MOCK EXAMINATION 3
% ==============================================================================
\section{University Mock Examination 3}

\begin{mockexam}[INT108 Mock Examination 3 (100 Marks --- 3 Hours)]
\noindent\textbf{Section A: Short Compulsory Questions ($10 \times 2 = 20\text{ Marks}$)}
\begin{enumerate}[leftmargin=1.5em]
    \item What is the difference between a syntax error and an exception?
    \item Explain the output of \texttt{print(bool([]), bool([0]))}.
    \item What is an accumulator variable in loop processing?
    \item Write the output of \texttt{"Python"[1:5:2]}.
    \item What is the purpose of the \texttt{nonlocal} keyword?
    \item Why are dictionary keys required to be hashable?
    \item What is the default file opening mode in Python?
    \item Differentiate between \texttt{@classmethod} and \texttt{@staticmethod}.
    \item What does \texttt{re.findall(r"\textbackslash S+", "Hi all")} return?
    \item Explain the meaning of the \texttt{\_\_name\_\_ == "\_\_main\_\_"} idiom.
\end{enumerate}

\medskip
\noindent\textbf{Section B: Analytical Medium Problems ($6 \times 5 = 30\text{ Marks}$)}
\begin{enumerate}[leftmargin=1.5em, start=11]
    \item Predict the output:
    \begin{lstlisting}[language=Python3]
s = "engineering"
d = {c: s.count(c) for c in set(s) if s.count(c) > 1}
print(sorted(d.items()))
    \end{lstlisting}
    \item Write a function to calculate the root-mean-square (RMS) of a numeric list.
    \item Implement a custom exception class \texttt{InvalidAgeError} raised when an age is outside $0-120$.
    \item Explain C3 Linearization and trace the MRO of multiple inheritance.
    \item Write a program to merge two sorted lists into a single sorted list in $O(N+M)$ time.
    \item Write a regex to validate an IPv4 address (e.g., \texttt{192.168.1.1}).
\end{enumerate}

\medskip
\noindent\textbf{Section C: Comprehensive Programming Problems ($5 \times 10 = 50\text{ Marks}$)}
\begin{enumerate}[leftmargin=1.5em, start=17]
    \item Implement a recursive function to compute the Tower of Hanoi for 4 disks, printing each step. (10 Marks)
    \item Design a \texttt{Polynomial} class representing $a x^2 + b x + c$ with an \texttt{add()} and \texttt{\_\_str\_\_} method. (10 Marks)
    \item Write a program to read a text file, count uppercase letters, lowercase letters, digits, and spaces, writing summary statistics to a new file. (10 Marks)
    \item Create a bank ATM withdrawal simulator with PIN retry limits (max 3 attempts) and transaction logs. (10 Marks)
    \item Write a function that uses regex to find and replace all phone numbers in a document with masked digits (\texttt{+91-XXXXX-XXXXX}). (10 Marks)
\end{enumerate}
\end{mockexam}

\subsection{Complete Worked Solutions for Mock Examination 3}

\begin{solutionbox}[Solutions to Mock Exam 3 (100 Marks)]
\noindent\textbf{Section A Solutions ($10 \times 2 = 20\text{ Marks}$)}:
\begin{enumerate}
    \item Syntax errors prevent parsing before execution; exceptions are runtime errors during execution. [2M]
    \item \texttt{bool([])} is \texttt{False} (empty); \texttt{bool([0])} is \texttt{True} (non-empty list). [2M]
    \item An accumulator variable aggregates values across loop iterations (e.g., \texttt{total += x}). [2M]
    \item \texttt{"yh"}. [2M]
    \item \texttt{nonlocal} binds identifiers to variables in outer enclosing functions. [2M]
    \item Dictionary keys use hash values for $O(1)$ table lookups; mutability would alter hash codes. [2M]
    \item Default mode is read text: \texttt{'r'}. [2M]
    \item \texttt{@classmethod} receives \texttt{cls}; \texttt{@staticmethod} receives neither \texttt{self} nor \texttt{cls}. [2M]
    \item \texttt{['Hi', 'all']} (extracts non-whitespace words). [2M]
    \item Ensures module code runs only when executed directly, not when imported. [2M]
\end{enumerate}

\medskip
\noindent\textbf{Section B Solutions ($6 \times 5 = 30\text{ Marks}$)}:
\begin{enumerate}[start=11]
    \item Character frequency dictionary for letters appearing $>1$ time. Output: \texttt{[('e', 3), ('g', 2), ('i', 2), ('n', 3)]}. [5M]
    \item RMS function:
    \begin{lstlisting}[language=Python3]
def rms(nums):
    return (sum(x**2 for x in nums) / len(nums)) ** 0.5
    \end{lstlisting} [5M]
    \item Custom exception:
    \begin{lstlisting}[language=Python3]
class InvalidAgeError(Exception): pass
def verify_age(age):
    if not 0 <= age <= 120: raise InvalidAgeError(f"Age {age} out of bounds")
    \end{lstlisting} [5M]
    \item MRO determines method lookup order in multiple inheritance using C3 Linearization. [5M]
    \item Linear Merge:
    \begin{lstlisting}[language=Python3]
def merge_sorted(l1, l2):
    i, j, res = 0, 0, []
    while i < len(l1) and j < len(l2):
        if l1[i] < l2[j]: res.append(l1[i]); i += 1
        else: res.append(l2[j]); j += 1
    return res + l1[i:] + l2[j:]
    \end{lstlisting}
    \item IPv4 regex: [5M]
    \begin{lstlisting}[language=Python3]
r"^(?:(?:25[0-5]|2[0-4]\d|[01]?\d\d?)\.){3}(?:25[0-5]|2[0-4]\d|[01]?\d\d?)$"
    \end{lstlisting}
\end{enumerate}

\medskip
\noindent\textbf{Section C Solutions ($5 \times 10 = 50\text{ Marks}$)}:
\begin{enumerate}[start=17]
    \item \textbf{Hanoi 4 Disks}:
    \begin{lstlisting}[language=Python3]
def hanoi(n, source, target, auxiliary):
    if n == 1:
        print(f"Move disk 1 from {source} -> {target}")
        return
    hanoi(n - 1, source, auxiliary, target)
    print(f"Move disk {n} from {source} -> {target}")
    hanoi(n - 1, auxiliary, target, source)
    \end{lstlisting}
    Requires $2^4 - 1 = 15$ moves recursively transferring sub-towers. [10M]
    \item \textbf{Polynomial Class}:
    \begin{lstlisting}[language=Python3]
class Polynomial:
    def __init__(self, a, b, c): self.a, self.b, self.c = a, b, c
    def add(self, o): return Polynomial(self.a+o.a, self.b+o.b, self.c+o.c)
    def __str__(self): return f"{self.a}x^2 + {self.b}x + {self.c}"
    \end{lstlisting} [10M]
    \item \textbf{File Stats Parser}:
    \begin{lstlisting}[language=Python3]
u = l = d = s = 0
with open("data.txt", "r") as f:
    for char in f.read():
        if char.isupper(): u += 1
        elif char.islower(): l += 1
        elif char.isdigit(): d += 1
        elif char.isspace(): s += 1
print(f"Upper: {u}, Lower: {l}, Digits: {d}, Spaces: {s}")
    \end{lstlisting} [10M]
    \item \textbf{ATM Simulator}:
    \begin{lstlisting}[language=Python3]
pin = "1234"; balance = 10000; attempts = 0
while attempts < 3:
    if input("Enter PIN: ") == pin:
        amt = float(input("Amount to withdraw: "))
        if amt <= balance: balance -= amt; print("Success! New Balance:", balance)
        else: print("Insufficient funds")
        break
    attempts += 1
    print(f"Wrong PIN! {3-attempts} attempts left.")
    \end{lstlisting} [10M]
    \item \textbf{Regex Masker}:
    \begin{lstlisting}[language=Python3]
import re
def mask_phones(doc):
    return re.sub(r"\+91-\d{5}-\d{5}", "+91-XXXXX-XXXXX", doc)
    \end{lstlisting} [10M]
\end{enumerate}
\end{solutionbox}

% ==============================================================================
% MOCK EXAMINATION 4
% ==============================================================================
\section{University Mock Examination 4}

\begin{mockexam}[INT108 Mock Examination 4 (100 Marks --- 3 Hours)]
\noindent\textbf{Section A: Short Compulsory Questions ($10 \times 2 = 20\text{ Marks}$)}
\begin{enumerate}[leftmargin=1.5em]
    \item What is the difference between interactive mode and script mode in Python?
    \item Evaluate: \texttt{print(10 > 5 and 3 < 1 or not 0)}.
    \item What is integer pooling / small integer caching in CPython?
    \item How do you create an empty set in Python?
    \item Explain the return value of \texttt{string.split(",")}.
    \item What is the purpose of the \texttt{\_\_str\_\_} dunder method?
    \item Differentiate between \texttt{os.mkdir()} and \texttt{os.makedirs()}.
    \item What is the difference between \texttt{pickle} and \texttt{json} modules?
    \item Explain the regex metacharacter \texttt{?} and its dual meaning.
    \item What is the output of \texttt{print(type(lambda x: x))}?
\end{enumerate}

\medskip
\noindent\textbf{Section B: Analytical Medium Problems ($6 \times 5 = 30\text{ Marks}$)}
\begin{enumerate}[leftmargin=1.5em, start=11]
    \item Predict the output of this dictionary comprehension:
    \begin{lstlisting}[language=Python3]
keys = ["a", "b", "c"]
vals = [1, 2, 3]
d = {k: v**2 for k, v in zip(keys, vals)}
print(d)
    \end{lstlisting}
    \item Write a function to check whether two input strings are anagrams of each other.
    \item Implement a generator function \texttt{even\_generator(limit)} using \texttt{yield}.
    \item Design a class \texttt{Circle} with a property decorator for \texttt{radius} and calculated \texttt{area}.
    \item Explain the execution behavior of \texttt{try-except-else-finally} when a return statement is inside \texttt{try}.
    \item Write a regex to extract all HTML anchor tags (\texttt{<a href="...">text</a>}).
\end{enumerate}

\medskip
\noindent\textbf{Section C: Comprehensive Programming Problems ($5 \times 10 = 50\text{ Marks}$)}
\begin{enumerate}[leftmargin=1.5em, start=17]
    \item Implement a complete Binary Search Tree (BST) class with \texttt{insert()} and \texttt{inorder\_traversal()} methods. (10 Marks)
    \item Write a Python program to perform Run-Length Encoding (RLE) string compression and decompression. (10 Marks)
    \item Design a multi-level inheritance hierarchy: \texttt{Vehicle} $\to$ \texttt{Car} $\to$ \texttt{ElectricCar} with battery management and range calculation. (10 Marks)
    \item Write a program to read a CSV gradebook file, compute course median and standard deviation, and write a summary report. (10 Marks)
    \item Implement a complete student attendance tracker with daily marking, percentage computation, and attendance shortage alerts ($<75\%$). (10 Marks)
\end{enumerate}
\end{mockexam}

\subsection{Complete Worked Solutions for Mock Examination 4}

\begin{solutionbox}[Solutions to Mock Exam 4 (100 Marks)]
\noindent\textbf{Section A Solutions ($10 \times 2 = 20\text{ Marks}$)}:
\begin{enumerate}
    \item Interactive mode evaluates expressions immediately line-by-line; script mode executes complete source files. [2M]
    \item \texttt{10 > 5 and 3 < 1} is \texttt{False}; \texttt{not 0} is \texttt{True}; \texttt{False or True} is \texttt{True}. [2M]
    \item CPython pre-caches integer objects in range $[-5, 256]$ to optimize memory allocation. [2M]
    \item \texttt{s = set()} (since \texttt{\{\}} creates an empty dictionary). [2M]
    \item Splits string by comma delimiter and returns a list of substring tokens. [2M]
    \item Provides user-friendly string representation when object is passed to \texttt{print()} or \texttt{str()}. [2M]
    \item \texttt{mkdir()} creates a single directory; \texttt{makedirs()} creates all intermediate parent directories recursively. [2M]
    \item \texttt{json} produces human-readable text; \texttt{pickle} produces Python-specific binary bytecode. [2M]
    \item \texttt{?} means 0 or 1 occurrences (optional), or makes quantifiers non-greedy. [2M]
    \item \texttt{<class 'function'>}. [2M]
\end{enumerate}

\medskip
\noindent\textbf{Section B Solutions ($6 \times 5 = 30\text{ Marks}$)}:
\begin{enumerate}[start=11]
    \item Dictionary comprehension with squared values: \texttt{\{'a': 1, 'b': 4, 'c': 9\}}. [5M]
    \item Anagram checker: \texttt{sorted(s1.lower()) == sorted(s2.lower())}. [5M]
    \item Generator:
    \begin{lstlisting}[language=Python3]
def even_gen(limit):
    for i in range(0, limit + 1, 2): yield i
    \end{lstlisting} [5M]
    \item Circle Property:
    \begin{lstlisting}[language=Python3]
import math
class Circle:
    def __init__(self, r): self._r = r
    @property
    def radius(self): return self._r
    @property
    def area(self): return math.pi * (self._r ** 2)
    \end{lstlisting} [5M]
    \item \texttt{finally} executes before the function exits, even if \texttt{return} is in \texttt{try}. [5M]
    \item Anchor regex: \lstinline|r'<a\s+href="[^"]*">(.*?)</a>'|. [5M]
\end{enumerate}

\medskip
\noindent\textbf{Section C Solutions ($5 \times 10 = 50\text{ Marks}$)}:
\begin{enumerate}[start=17]
    \item \textbf{BST Implementation}:
    \begin{lstlisting}[language=Python3]
class Node:
    def __init__(self, val): self.val = val; self.left = self.right = None

class BST:
    def __init__(self): self.root = None
    def insert(self, val):
        if not self.root: self.root = Node(val); return
        curr = self.root
        while True:
            if val < curr.val:
                if not curr.left: curr.left = Node(val); break
                curr = curr.left
            else:
                if not curr.right: curr.right = Node(val); break
                curr = curr.right
    def inorder_traversal(self, node=None, is_root=True):
        if is_root: node = self.root
        if node:
            self.inorder_traversal(node.left, False)
            print(node.val, end=" ")
            self.inorder_traversal(node.right, False)
    \end{lstlisting} [10M]
    \item \textbf{RLE Compression}:
    \begin{lstlisting}[language=Python3]
def rle_compress(text):
    if not text: return ""
    res, count = [], 1
    for i in range(1, len(text)):
        if text[i] == text[i-1]: count += 1
        else: res.append(f"{text[i-1]}{count}"); count = 1
    res.append(f"{text[-1]}{count}")
    return "".join(res)
# Tests: rle_compress("AABBB") -> "A2B3", rle_compress("X") -> "X1"
    \end{lstlisting} [10M]
    \item \textbf{Electric Car Hierarchy}:
    \begin{lstlisting}[language=Python3]
class Vehicle:
    def __init__(self, brand): self.brand = brand
class Car(Vehicle):
    def __init__(self, brand, model):
        super().__init__(brand); self.model = model
class ElectricCar(Car):
    def __init__(self, brand, model, battery, efficiency):
        super().__init__(brand, model)
        self.battery = battery
        self.efficiency = efficiency
    def range(self): return self.battery / self.efficiency
# Test: ev = ElectricCar("Tesla", "M3", 75, 0.15); ev.range() -> 500.0
    \end{lstlisting} [10M]
    \item \textbf{CSV Statistics Processor}:
    \begin{lstlisting}[language=Python3]
import csv
def process_scores(filename):
    scores = []
    with open(filename, "r") as f:
        reader = csv.reader(f)
        next(reader) # Skip header
        for row in reader: scores.append(float(row[1]))
    n = len(scores)
    mean = sum(scores) / n
    scores.sort()
    med = scores[n//2] if n % 2 != 0 else (scores[n//2-1]+scores[n//2])/2
    var = sum((x - mean)**2 for x in scores) / n
    return mean, med, var
    \end{lstlisting} [10M]
    \item \textbf{Attendance Tracker}:
    \begin{lstlisting}[language=Python3]
class AttendanceTracker:
    def __init__(self): self.records = {}
    def mark(self, reg, present=True):
        p, t = self.records.get(reg, (0, 0))
        self.records[reg] = (p + (1 if present else 0), t + 1)
    def check_shortages(self):
        for reg, (p, t) in self.records.items():
            pct = (p / t) * 100
            if pct < 75.0: print(f"ALERT: {reg} has {pct:.1f}% attendance (Shortage!)")
    \end{lstlisting} [10M]
\end{enumerate}
\end{solutionbox}

% ==============================================================================
% MOCK EXAMINATION 5
% ==============================================================================
\section{University Mock Examination 5}

\begin{mockexam}[INT108 Mock Examination 5 (100 Marks --- 3 Hours)]
\noindent\textbf{Section A: Short Compulsory Questions ($10 \times 2 = 20\text{ Marks}$)}
\begin{enumerate}[leftmargin=1.5em]
    \item Explain the difference between mutable and immutable objects with two examples of each.
    \item What is the output of \texttt{print("abc".center(10, "*"))}?
    \item State the purpose of the \texttt{assert} statement in Python debugging.
    \item What is the difference between \texttt{list.append(x)} and \texttt{list.insert(0, x)} in terms of time complexity?
    \item How does Python resolve names in nested functions?
    \item Explain the role of \texttt{\_\_del\_\_} destructor in Python.
    \item What does \texttt{re.search(r"\textasciicircum\textbackslash d\{3\}-\textbackslash d\{3\}\$", "123-456")} return?
    \item State the purpose of the \texttt{math.isclose()} function.
    \item How can you open a file for both reading and writing without truncating its content?
    \item What is the difference between \texttt{range(10)} and \texttt{list(range(10))} in memory?
\end{enumerate}

\medskip
\noindent\textbf{Section B: Analytical Medium Problems ($6 \times 5 = 30\text{ Marks}$)}
\begin{enumerate}[leftmargin=1.5em, start=11]
    \item Predict the output of this inheritance and class variable tracing program:
    \begin{lstlisting}[language=Python3]
class Parent:
    val = 10
class Child1(Parent):
    pass
class Child2(Parent):
    pass
Child1.val = 20
Parent.val = 30
print(Parent.val, Child1.val, Child2.val)
    \end{lstlisting}
    \item Write a function to check if a matrix is an identity matrix ($I\_n$).
    \item Implement a generator function \texttt{fibonacci\_gen(n)} that yields the first $n$ Fibonacci numbers.
    \item Write a program to count vowels, consonants, digits, and special characters in an input string.
    \item Explain how the \texttt{try-except-else-finally} block behaves when an uncaught exception is raised in \texttt{finally}.
    \item Write a regular expression to validate a strong password (min 8 chars, 1 uppercase, 1 lowercase, 1 digit, 1 special symbol).
\end{enumerate}

\medskip
\noindent\textbf{Section C: Comprehensive Programming Problems ($5 \times 10 = 50\text{ Marks}$)}
\begin{enumerate}[leftmargin=1.5em, start=17]
    \item Implement a complete \texttt{Queue} data structure class using a Python list with \texttt{enqueue()}, \texttt{dequeue()}, \texttt{is\_empty()}, and \texttt{size()} methods, raising custom exceptions on underflow. (10 Marks)
    \item Write a recursive function to compute the determinant of an $N \times N$ matrix. (10 Marks)
    \item Design an employee management system with an abstract \texttt{Employee} base class and subclasses \texttt{Manager}, \texttt{Developer}, and \texttt{Intern} calculating bonus payouts polymorphically. (10 Marks)
    \item Write a program to read a text file, find all palindromic words of length $\ge 3$, and write them in alphabetical order to \texttt{palindromes.txt}. (10 Marks)
    \item Build an automated log analyzer parsing Apache access logs with regex to compute total bandwidth transferred and top 5 client IPs. (10 Marks)
\end{enumerate}
\end{mockexam}

\subsection{Complete Worked Solutions for Mock Examination 5}

\begin{solutionbox}[Solutions to Mock Exam 5 (100 Marks)]
\noindent\textbf{Section A Solutions ($10 \times 2 = 20\text{ Marks}$)}:
\begin{enumerate}
    \item Mutable objects can be modified in place (\texttt{list}, \texttt{dict}); immutable objects cannot be altered after creation (\texttt{int}, \texttt{str}, \texttt{tuple}). [2M]
    \item \texttt{"***abc****"}. [2M]
    \item \texttt{assert condition, message} tests internal invariants, raising \texttt{AssertionError} if false. [2M]
    \item \texttt{append()} is $O(1)$ amortized; \texttt{insert(0, x)} is $O(n)$ due to shifting all existing elements. [2M]
    \item Via the LEGB rule, searching Local, then Enclosing outer function frames, then Global, then Built-in. [2M]
    \item \texttt{\_\_del\_\_} is called when an object's reference count drops to zero before garbage collection. [2M]
    \item Returns a \texttt{re.Match} object spanning \texttt{(0, 7)}. [2M]
    \item Tests whether two floating-point values are approximately equal within a relative tolerance to avoid float rounding errors. [2M]
    \item Using mode \texttt{'r+'}. [2M]
    \item \texttt{range(10)} is an $O(1)$ memory generator object; \texttt{list(range(10))} creates an $O(n)$ in-memory list containing 10 integers. [2M]
\end{enumerate}

\medskip
\noindent\textbf{Section B Solutions ($6 \times 5 = 30\text{ Marks}$)}:
\begin{enumerate}[start=11]
    \item Class variable inheritance tracing: \texttt{Child1.val = 20} sets attribute on \texttt{Child1}; \texttt{Parent.val = 30} updates \texttt{Parent} and \texttt{Child2} (which inherits from \texttt{Parent}). Output: \texttt{30 20 30}. [5M]
    \item Identity matrix check:
    \begin{lstlisting}[language=Python3]
def is_identity(m):
    n = len(m)
    return all(m[i][j] == (1 if i == j else 0) for i in range(n) for j in range(n))
    \end{lstlisting} [5M]
    \item Fibonacci generator:
    \begin{lstlisting}[language=Python3]
def fib_gen(n):
    a, b = 0, 1
    for _ in range(n):
        yield a
        a, b = b, a + b
    \end{lstlisting} [5M]
    \item Character classification:
    \begin{lstlisting}[language=Python3]
def count_chars(text):
    v = c = d = s = 0
    for ch in text:
        if ch.lower() in "aeiou": v += 1
        elif ch.isalpha(): c += 1
        elif ch.isdigit(): d += 1
        else: s += 1
    return v, c, d, s
    \end{lstlisting} [5M]
    \item An exception in \texttt{finally} supersedes any previous exceptions or return values, propagating immediately to the caller. [5M]
    \item Password regex: \texttt{r"\textasciicircum(?=.*[A-Z])(?=.*[a-z])(?=.*\textbackslash d)(?=.*[!@\#\$\%\textasciicircum\&\*]).\{8,\}\$"}. [5M]
\end{enumerate}

\medskip
\noindent\textbf{Section C Solutions ($5 \times 10 = 50\text{ Marks}$)}:
\begin{enumerate}[start=17]
    \item \textbf{Queue Class with Custom Exception}:
    \begin{lstlisting}[language=Python3]
class QueueUnderflowError(Exception): pass
class Queue:
    def __init__(self): self.items = []
    def enqueue(self, item): self.items.append(item)
    def dequeue(self):
        if not self.items: raise QueueUnderflowError("Queue is empty")
        return self.items.pop(0)
    def is_empty(self): return len(self.items) == 0
    def size(self): return len(self.items)
    \end{lstlisting} [10M]
    \item \textbf{Recursive Determinant}:
    \begin{lstlisting}[language=Python3]
def det(m):
    n = len(m)
    if n == 1: return m[0][0]
    if n == 2: return m[0][0]*m[1][1] - m[0][1]*m[1][0]
    total = 0
    for c in range(n):
        sub = [row[:c] + row[c+1:] for row in m[1:]]
        total += ((-1)**c) * m[0][c] * det(sub)
    return total
    \end{lstlisting} [10M]
    \item \textbf{Polymorphic Employee Bonus System}:
    \begin{lstlisting}[language=Python3]
class Employee:
    def __init__(self, name, salary):
        self.name = name; self.salary = salary
    def get_bonus(self): raise NotImplementedError

class Manager(Employee):
    def get_bonus(self): return self.salary * 0.20

class Developer(Employee):
    def get_bonus(self): return self.salary * 0.15

class Intern(Employee):
    def get_bonus(self): return self.salary * 0.05

# Tests
mgr = Manager("Alice", 100000)
dev = Developer("Bob", 80000)
print(mgr.get_bonus(), dev.get_bonus()) # 20000.0, 12000.0
    \end{lstlisting} [10M]
    \item \textbf{Palindromic Words Parser}:
    \begin{lstlisting}[language=Python3]
palindromes = set()
with open("input.txt", "r") as f:
    for word in f.read().lower().split():
        clean = "".join(ch for ch in word if ch.isalnum())
        if len(clean) >= 3 and clean == clean[::-1]:
            palindromes.add(clean)
with open("palindromes.txt", "w") as out:
    for p in sorted(palindromes): out.write(p + "\n")
    \end{lstlisting} [10M]
    \item \textbf{Log Bandwidth \& Top IPs}:
    \begin{lstlisting}[language=Python3]
import re
from collections import defaultdict

def analyze_logs(logfile):
    ip_bytes = defaultdict(int)
    total_bandwidth = 0
    pattern = r'(\d+\.\d+\.\d+\.\d+).*" \d+ (\d+)'
    
    with open(logfile, "r") as f:
        for line in f:
            match = re.search(pattern, line)
            if match:
                ip, bytes_sent = match.groups()
                b = int(bytes_sent)
                ip_bytes[ip] += b
                total_bandwidth += b
                
    top_ips = sorted(ip_bytes.items(), key=lambda x: x[1], reverse=True)[:5]
    print(f"Total Bandwidth: {total_bandwidth} bytes")
    for ip, b in top_ips: print(f"{ip}: {b} bytes")
    \end{lstlisting} [10M]
\end{enumerate}
\end{solutionbox}

% ==============================================================================
% MOCK EXAMINATION 6
% ==============================================================================
\section{University Mock Examination 6}

\begin{mockexam}[INT108 Mock Examination 6 (100 Marks --- 3 Hours)]
\noindent\textbf{Section A: Short Compulsory Questions ($10 \times 2 = 20\text{ Marks}$)}
\begin{enumerate}[leftmargin=1.5em]
    \item Explain the difference between \texttt{break} in a \texttt{for} loop vs. \texttt{sys.exit()}.
    \item What is the value of \texttt{3 << 3} and \texttt{24 >> 3}?
    \item How do you check if a key exists in a dictionary without raising \texttt{KeyError}?
    \item What is the output of \texttt{list("hello")} and \texttt{"".join(['h', 'i'])}?
    \item Explain why \texttt{isinstance(True, int)} evaluates to \texttt{True}.
    \item What is the difference between \texttt{f.seek(0)} and \texttt{f.seek(0, 2)}?
    \item What is an anonymous lambda function and what are its syntax limits?
    \item Differentiate between \texttt{re.findall()} and \texttt{re.finditer()}.
    \item What does \texttt{round(2.5)} vs \texttt{round(3.5)} return under Python's bankers' rounding?
    \item Explain the term duck typing in Python.
\end{enumerate}

\medskip
\noindent\textbf{Section B: Analytical Medium Problems ($6 \times 5 = 30\text{ Marks}$)}
\begin{enumerate}[leftmargin=1.5em, start=11]
    \item Predict the output:
    \begin{lstlisting}[language=Python3]
def outer():
    x = 10
    def inner():
        nonlocal x
        x += 5
        return x
    return inner

f1 = outer()
print(f1(), f1())
    \end{lstlisting}
    \item Write a function to check if a number is a Strong Number ($n = \sum \text{digit}!$).
    \item Implement a matrix rotation algorithm that rotates an $N \times N$ matrix 90 degrees clockwise in-place.
    \item Write a Python context manager class \texttt{Timer} that measures code execution time using \texttt{\_\_enter\_\_} and \texttt{\_\_exit\_\_}.
    \item Explain the mutable default argument trap and write the standard idiom to prevent it.
    \item Write a regex to extract all valid dates in format \texttt{DD-MM-YYYY} with day and month bounds checking.
\end{enumerate}

\medskip
\noindent\textbf{Section C: Comprehensive Programming Problems ($5 \times 10 = 50\text{ Marks}$)}
\begin{enumerate}[leftmargin=1.5em, start=17]
    \item Implement a complete \texttt{Stack} data structure with \texttt{push()}, \texttt{pop()}, \texttt{peek()}, and min-element retrieval in $O(1)$ time. (10 Marks)
    \item Write a program to solve the Sudoku 1D row/column/box validation check for a $9 \times 9$ grid. (10 Marks)
    \item Design a complete banking account hierarchy with savings accounts, current accounts, overdraft limits, and interest accrual. (10 Marks)
    \item Write a program that compresses text using Run-Length Encoding and writes both compressed and decompressed verification files. (10 Marks)
    \item Build a complete course grade management system that reads CSV records, calculates weighted averages (Assignments 30\%, Midterm 30\%, Final 40\%), and generates formatted letter grades. (10 Marks)
\end{enumerate}
\end{mockexam}

\subsection{Complete Worked Solutions for Mock Examination 6}

\begin{solutionbox}[Solutions to Mock Exam 6 (100 Marks)]
\noindent\textbf{Section A Solutions ($10 \times 2 = 20\text{ Marks}$)}:
\begin{enumerate}
    \item \texttt{break} exits only the innermost loop; \texttt{sys.exit()} terminates the entire Python process. [2M]
    \item $3 \times 2^3 = \mathbf{24}$; $24 // 2^3 = \mathbf{3}$. [2M]
    \item Using \texttt{key in dict} or \texttt{dict.get(key, default)}. [2M]
    \item \texttt{['h', 'e', 'l', 'l', 'o']} and \texttt{"hi"}. [2M]
    \item \texttt{bool} is a subclass of \texttt{int} in Python (\texttt{True == 1}, \texttt{False == 0}). [2M]
    \item \texttt{f.seek(0)} moves pointer to start; \texttt{f.seek(0, 2)} moves pointer to end of file. [2M]
    \item A single-expression inline function; it cannot contain multiple statements, loops, or assignments. [2M]
    \item \texttt{findall()} returns a list of strings; \texttt{finditer()} yields an iterator of \texttt{Match} objects with start/end indices. [2M]
    \item Python uses round-to-even: \texttt{round(2.5) = 2} and \texttt{round(3.5) = 4}. [2M]
    \item "If it walks like a duck and quacks like a duck, it's a duck" — object validity is determined by the presence of methods, not explicit inheritance. [2M]
\end{enumerate}

\medskip
\noindent\textbf{Section B Solutions ($6 \times 5 = 30\text{ Marks}$)}:
\begin{enumerate}[start=11]
    \item Closure state persistence: first call returns $10+5 = 15$; second call returns $15+5 = 20$. Output: \texttt{15 20}. [5M]
    \item Strong number checker:
    \begin{lstlisting}[language=Python3]
import math
def is_strong(n):
    return sum(math.factorial(int(d)) for d in str(n)) == n
    \end{lstlisting} [5M]
    \item 90-degree Matrix Rotation:
    \begin{lstlisting}[language=Python3]
def rotate_90(mat):
    # Transpose then reverse rows
    n = len(mat)
    for i in range(n):
        for j in range(i, n): mat[i][j], mat[j][i] = mat[j][i], mat[i][j]
    for row in mat: row.reverse()
    \end{lstlisting} [5M]
    \item Timer Context Manager:
    \begin{lstlisting}[language=Python3]
import time
class Timer:
    def __enter__(self): self.start = time.time(); return self
    def __exit__(self, *args): print(f"Elapsed: {time.time() - self.start:.4f}s")
    \end{lstlisting} [5M]
    \item Default evaluated once. Standard fix: \texttt{def func(val, item=None): if item is None: item = []}. [5M]
    \item Date regex: \texttt{r"\textasciicircum(0[1-9]|[12]\textbackslash d|3[01])-(0[1-9]|1[0-2])-\textbackslash d\{4\}\$"}. [5M]
\end{enumerate}

\medskip
\noindent\textbf{Section C Solutions ($5 \times 10 = 50\text{ Marks}$)}:
\begin{enumerate}[start=17]
    \item \textbf{Min-Stack in O(1)}:
    \begin{lstlisting}[language=Python3]
class MinStack:
    def __init__(self): self.stack = []; self.min_stack = []
    def push(self, x):
        self.stack.append(x)
        if not self.min_stack or x <= self.min_stack[-1]: self.min_stack.append(x)
    def pop(self):
        val = self.stack.pop()
        if val == self.min_stack[-1]: self.min_stack.pop()
        return val
    def get_min(self): return self.min_stack[-1]
    \end{lstlisting} [10M]
    \item \textbf{Sudoku Validation}:
    \begin{lstlisting}[language=Python3]
def is_valid_sudoku(board):
    def is_valid(group):
        nums = [n for n in group if n != '.']
        return len(nums) == len(set(nums))
    for row in board:
        if not is_valid(row): return False
    for col in zip(*board):
        if not is_valid(col): return False
    for i in (0, 3, 6):
        for j in (0, 3, 6):
            box = [board[x][y] for x in range(i, i+3) for y in range(j, j+3)]
            if not is_valid(box): return False
    return True
# Test: pass a 9x9 list of lists of strings
    \end{lstlisting} [10M]
    \item \textbf{Banking Account Hierarchy}:
    \begin{lstlisting}[language=Python3]
class Account:
    def __init__(self, acc_id, bal): self.acc_id, self.bal = acc_id, bal
    def deposit(self, amt): self.bal += amt
    def withdraw(self, amt):
        if amt <= self.bal: self.bal -= amt; return True
        return False
class SavingsAccount(Account):
    def __init__(self, acc_id, bal, rate=0.03):
        super().__init__(acc_id, bal); self.rate = rate
    def add_interest(self): self.deposit(self.bal * self.rate)
class CurrentAccount(Account):
    def __init__(self, acc_id, bal, limit=5000):
        super().__init__(acc_id, bal); self.limit = limit
    def withdraw(self, amt):
        if self.bal + self.limit >= amt: self.bal -= amt; return True
        return False
    \end{lstlisting} [10M]
    \item \textbf{RLE Compression \& Verification}:
    \begin{lstlisting}[language=Python3]
def rle_file_compress(infile, outfile):
    with open(infile, "r") as f: data = f.read()
    if not data: return
    comp, count = [], 1
    for i in range(1, len(data)):
        if data[i] == data[i-1]: count += 1
        else: comp.append(f"{data[i-1]}{count}"); count = 1
    comp.append(f"{data[-1]}{count}")
    with open(outfile, "w") as f: f.write("".join(comp))

def rle_verify(infile, compfile):
    with open(infile, "r") as f: orig = f.read()
    with open(compfile, "r") as f: comp = f.read()
    import re
    decomp = "".join(char * int(num) for char, num in re.findall(r'(.)(\d+)', comp))
    assert orig == decomp, "Mismatch!"
    print("Verification Passed")
    \end{lstlisting} [10M]
    \item \textbf{Weighted Grade Calculator}:
    \begin{lstlisting}[language=Python3]
import csv
def calculate_grades(infile, outfile):
    with open(infile, "r") as fin, open(outfile, "w") as fout:
        reader = csv.reader(fin); writer = csv.writer(fout)
        header = next(reader); writer.writerow(header + ["Final", "Grade"])
        for row in reader:
            assign, mid, final = float(row[1]), float(row[2]), float(row[3])
            score = 0.3*assign + 0.3*mid + 0.4*final
            grade = "A" if score >= 90 else "B" if score >= 80 else "C" if score >= 70 else "F"
            writer.writerow(row + [f"{score:.1f}", grade])
# Test Data: Student1, 85, 90, 88 -> Score: 87.7, Grade: B
    \end{lstlisting} [10M]
\end{enumerate}
\end{solutionbox}
